% Adapted from Wenxuan Shang's MIT-licensed NKU lab report template.
% Original license retained in TEMPLATE-LICENSE.
\documentclass[UTF8,notitlepage,11pt]{ctexart}
\usepackage[a4paper,margin=1in]{geometry}
\usepackage{amsmath,graphicx,booktabs,tabularx,longtable,array,float}
\usepackage{fancyhdr,fvextra,xcolor,enumitem,tikz}
\usetikzlibrary{arrows.meta,positioning,fit,backgrounds}
\usepackage[colorlinks=true,linkcolor=blue!45!black,urlcolor=blue!55!black]{hyperref}
\usepackage{xurl}
\setlength{\headheight}{15pt}\pagestyle{fancy}\fancyhf{}
\fancyhead[L]{尚文轩 / 2412456}\fancyhead[R]{计算机网络 · 实验二}\fancyfoot[C]{\thepage}
\setlength{\parskip}{4pt}\renewcommand{\arraystretch}{1.2}
\setlist{nosep,topsep=4pt}
\DefineVerbatimEnvironment{code}{Verbatim}{fontsize=\small,breaklines=true,frame=single}
\newcommand{\src}[1]{\nolinkurl{#1}}
\newcommand{\ParserTests}{24}
\newcommand{\IntegrationTests}{11}
\newcommand{\TestDuration}{6.797}
\newcommand{\TestTimestamp}{2026-10-22 00:37:12 +0800}
\newcommand{\PacketCount}{6}
\newcommand{\CaptureDuration}{4.012}
\newcommand{\DriverReceived}{6}
\newcommand{\DriverDropped}{0}
\newcommand{\InterfaceDropped}{0}
\newcommand{\ShortFrames}{0}
\newcommand{\TruncatedFrames}{0}

\title{\vspace*{-1.2cm}\includegraphics[width=.95\textwidth]{nankai.jpg}\\[1.2em]
{\bfseries\Huge 实验报告}\\[1em]{\LARGE 实验02：编写数据包捕获程序\\并对捕获的数据包进行分析}}
\author{\begin{minipage}{.85\textwidth}\centering\Large
姓名：\underline{尚文轩}\\[.6em]
学号：\underline{2412456}\\[.6em]
专业：\underline{计算机科学与技术}\\[.6em]
任课教师：\underline{徐敬东、张建忠}\end{minipage}}
\date{\Large 2026年10月22日}
\begin{document}
\maketitle\thispagestyle{empty}\clearpage
\pagenumbering{arabic}
{\small\setlength{\parskip}{0pt}\tableofcontents}
\clearpage

\section{实验目的与要求对应}
本实验使用 C++17 与 Npcap，在 Windows 上枚举网卡、打开指定设备，捕获本机以太网视图下的数据帧，并在屏幕上显示源 MAC 地址、目的 MAC 地址、类型/长度字段与该字段之后第一个字节。
通过对真实 ICMP 请求与应答帧的逐字节复核，理解链路层封装、网络字节序和捕获长度的含义；通过独立边界测试验证解析器不会在截断数据上越界。

\begin{table}[H]\centering\small
\caption{任务要求、实现与验证证据}
\begin{tabularx}{\textwidth}{p{.25\textwidth}XX}\toprule
要求 & 实现内容 & 证据\\\midrule
了解 Npcap 架构 & 用户态 API、DLL、NDIS 过滤驱动、适配器分层 & 第\ref{sec:architecture}节架构图\\
设备枚举与打开 & findalldevs、create、配置、activate & 设备列表与真实网卡测试\\
捕获并显示必需字段 & next\_ex、显式按字节解析以太网头 & 真实 PCAP 与屏幕输出逐字段一致\\
简明屏幕显示，无需保存 & 主表四个必显字段，默认无落盘 & 默认不写抓包文件测试\\
结构清晰，可读性好 & 解析、Npcap 资源、命令行与捕获流程分模块 & 模块表、核心代码、异常清理\\
源码、可执行程序与报告 & C++、Windows x64 exe、LaTeX 与 PDF & 独立实验目录和白名单提交包\\\bottomrule
\end{tabularx}\end{table}

\section{实验环境与 Npcap 架构}\label{sec:architecture}
\subsection{实际环境}
操作系统为 Windows 11（10.0.26200）；编译器为 MSVC 19.29 x64，编译选项包含
\src{/std:c++17 /W4 /WX /utf-8 /O2 /MT}，警告作为错误处理，静态链接 C/C++ 运行库。
通过官方 Npcap SDK 1.16 的头文件获得 API 声明，运行时从系统 Npcap 目录加载 DLL。
实际捕获环境如下表，数据来自本机实测记录。

\begin{table}[H]\centering\small\caption{真实捕获的运行环境}
\begin{tabularx}{\textwidth}{lX}\toprule
Npcap & Npcap version 1.89, based on libpcap version 1.10.7 (64-bit time\_t) \\
网卡 & MediaTek Wi-Fi 6 MT7921 Wireless LAN Card \\
枚举编号 & 6（仅对本次枚举有效） \\
本机 IPv4 & 10.130.33.54 \\
默认网关 & 10.130.0.1 \\
过滤器 & icmp and host 10.130.0.1 \\
\bottomrule\end{tabularx}

\end{table}

本机原有旧 Npcap 驱动，但缺少配套 DLL；第一次安装失败，第二次由用户确认官方安装窗口完成安装。
后续以 DLL 实际版本、成功枚举网卡、打开设备及真实捕获验证环境，未将“下载完成”或“安装器启动”视为安装成功。
Npcap 安装时启用了管理员访问限制，因此真实捕获测试在用户确认的管理员进程中运行。

\subsection{架构与数据路径}
Npcap 提供 libpcap 风格的用户态 API，底层使用 Windows 内核驱动获取数据。其驱动作为 NDIS 6 轻量级过滤驱动连接网络栈和适配器；应用程序不直接读取网卡硬件寄存器，也不能用普通 TCP \src{recv} 代替链路层抓包\cite{npcapdev}。

\begin{figure}[H]\centering
\begin{tikzpicture}[>=Stealth,node distance=7mm,font=\small,
 box/.style={draw,rounded corners,text width=85mm,align=center,minimum height=10mm}]
\node[box] (app) {本实验 C++ 应用程序\\枚举、配置、捕获、解析与显示};
\node[box,below=of app] (lib) {wpcap.dll：libpcap API\\pcap\_findalldevs / pcap\_create / pcap\_next\_ex};
\node[box,below=of lib] (packet) {Packet.dll 与设备接口\\用户态调用进入捕获驱动};
\node[box,below=13mm of packet] (driver) {npcap.sys：NDIS 6 LWF\\过滤、内核缓冲与数据复制};
\node[box,below=of driver] (nic) {NDIS 网络栈与实际网卡\\本实验：Wi-Fi 普通模式的 Ethernet 视图};
\draw[<->] (app)--(lib);\draw[<->] (lib)--(packet);
\draw[<->] (packet)--node[right]{用户态 / 内核态}(driver);
\draw[<->] (driver)--(nic);
\end{tikzpicture}\caption{本实验使用的 Npcap 捕获数据路径（概念图）}
\end{figure}

\subsection{链路类型不能混用}
\src{pcap_datalink} 给出本句柄数据的链路类型。普通 Ethernet 或 Wi-Fi 适配器在本实验中返回
\src{DLT_EN10MB=1}，可以按以太网头解析；这不意味着抓到了原始 802.11 无线管理帧。
未启用无线监控模式，不把无线信道、重试标志或 Radiotap 等当成已捕获内容。

Npcap 回环通过 Windows Filtering Platform 支持本机流量，通常使用 \src{DLT_NULL}，只有地址族等伪首部，没有以太网源/目的 MAC\cite{npcapdev}。
因此程序对非 DLT\_EN10MB 句柄直接提示错误并退出，不将回环 IP 数据的前 12 字节伪装为 MAC。
离线 DLT\_NULL 合成文件也用于验证该拒绝路径。

\section{数据帧格式与解析规则}\label{sec:frame}
\subsection{以太网必显字段}
对 DLT\_EN10MB，以下偏移均从 Npcap 返回的数据缓冲起点计数，不包含前导码和 SFD。
MAC 按原始字节顺序输出，类型字段按大端解释。

\begin{table}[H]\centering
\caption{题目要求的字段与字节位置}
\begin{tabular}{cllc}\toprule
字节偏移 & 长度 & 含义 & 读取条件\\\midrule
0--5 & 6 & 目的 MAC & caplen $\ge 14$\\
6--11 & 6 & 源 MAC & caplen $\ge 14$\\
12--13 & 2 & 类型/长度字段 & caplen $\ge 14$\\
14 & 1 & 字段之后第一个字节 & caplen $\ge 15$\\\bottomrule
\end{tabular}\end{table}

程序把是否存在一个字段与字段的数值分开表示。例如 ARP 后首字节可能为 \src{0x00}，这仍是实际捕获到的有效数值，不能因为它等于零就显示 N/A。
若捕获长度仅为 14，则 MAC 与类型字段有效，但首字节不存在；若不足 14，则整条标记为短以太网帧。

\subsection{类型字段与长度字段的区分}
IEEE 802.3 长度字段的最大值是 1500；字段值达到 1536（\src{0x0600}）及以上才按 EtherType 解释；1501--1535 是保留区间。
本实验常见 EtherType 为 IPv4（\src{0x0800}）、ARP（\src{0x0806}）、IPv6（\src{0x86DD}）。
不能将所有 16 位数值都称为 IPv4 类型，也不能把某一条样本的首字节固定写入程序。类型数值参照 IEEE 802 Numbers 登记表\cite{ethertypes}。

\subsection{VLAN 的两个不同“首字节”}
对于 \src{0x8100} 或 \src{0x88A8}，偏移 14 开始是 VLAN 控制字段 TCI，之后才是内层 EtherType 和实际负载。
题目要求的“类型字段之后第一个字节”始终显示\textbf{外层字段之后的原始偏移 14}；程序另外显示内层类型、VLAN ID、负载偏移及负载首字节，避免混淆。

例如合成测试字节 \src{81 00 A0 64 08 00 45} 表示外层类型 \src{0x8100}，After 为 \src{0xA0}，TCI 为 \src{0xA064}，VLAN ID 为 100，内层类型为 IPv4，真正 IP 首字节为 \src{0x45}，负载从偏移 18 开始。双标签时继续向后解析；每次至少检查剩余四个字节，循环次数受实际捕获长度限制。
此处 VLAN 是明确标记的\textbf{合成边界测试}，本次真实网卡样本并未捕获 VLAN 标签。

\subsection{截断、字节序与解析范围}
\src{pcap_pkthdr.caplen} 是缓冲区实际可读取的字节数，\src{len} 是原始帧长度。
若 caplen 小于 len，只能在 caplen 范围内读取，并标注 TRUNCATED。不能因为原始帧足够长就读取未捕获的内存。
本程序不强转为带平台填充的头部结构体，使用显式字节操作；因此未对齐缓冲区也能安全解析。

\begin{code}
inline uint16_t read_be16(const uint8_t* p) {
    return static_cast<uint16_t>(
        (static_cast<uint16_t>(p[0]) << 8) | p[1]);
}
// MAC destination: data[0..5]; source: data[6..11].
// Type/length: read_be16(data + 12).
// Mandatory first byte: data[14], only if caplen >= 15.
\end{code}

\section{程序设计与实现}
\subsection{模块划分}
\begin{table}[H]\centering\small
\begin{tabularx}{\textwidth}{p{.30\textwidth}X}\toprule
文件 / 模块 & 职责\\\midrule
src/ethernet.hpp & 无驱动依赖的帧解析、MAC 格式化、类型分类、VLAN 与可选 IP/ARP 解释\\
src/npcap\_api.hpp & 系统 Npcap DLL 加载、SDK 函数类型解析、设备列表/句柄/转储资源回收\\
src/main.cpp & 参数校验、设备枚举、捕获流程、屏幕或 JSON 输出、退出和统计\\
tests/parser\_tests.cpp & 24 项纯 C++ 边界测试，无需 Npcap 运行库\\
tests/integration.py & 驱动真实 exe，比较捕获与 PCAP 原始字节\\
tests/run-elevated.ps1 & 选择物理默认路由网卡，记录环境，在管理员进程运行测试\\
tests/prepare\_report.py & 从真实保存证据生成报告表格和十六进制记录\\\bottomrule
\end{tabularx}\end{table}

\subsection{API 调用顺序与含义}
先使用官方 SDK 的函数声明，通过 \src{LoadLibraryExW} 加载系统 Npcap 目录的 DLL，并用 \src{GetProcAddress} 解析函数。
搜索目录限制为该 DLL 所在目录与系统目录，避免误用旧 WinPcap。
\src{pcap_init(PCAP_CHAR_ENC_UTF_8)} 明确约定 Npcap 参数、返回字符串的编码；程序 \src{wmain} 将 Windows UTF-16 参数转换成 UTF-8，中文路径的保存文件复核已通过测试\cite{pcapinit}。

\begin{table}[H]\centering\small
\caption{捕获生命周期中的关键 API}
\begin{tabularx}{\textwidth}{p{.39\textwidth}X}\toprule
API & 用途及检查\\\midrule
pcap\_findalldevs / freealldevs & 获取并释放设备链表；遍历名称、描述、IPv4 与标志\\
pcap\_create & 创建未激活句柄，失败检查错误缓冲\\
pcap\_set\_snaplen & 设置每帧捕获上限，默认 65535\\
pcap\_set\_promisc & 默认关闭混杂，显式选项可开启\\
pcap\_set\_timeout / buffer\_size & 设置 100 ms 缓冲读取期限和 2 MiB 捕获缓冲\\
pcap\_activate & 激活设备；负数为错误，正数警告单独显示\\
pcap\_datalink & 仅允许 DLT\_EN10MB 的以太网解析路径\\
pcap\_compile / setfilter / freecode & 可选 BPF 编译、安装与释放；过滤失败不能默默继续\\
pcap\_setnonblock & 设为非阻塞，确保空闲时仍能检测停止条件\\
pcap\_next\_ex & 读取下一帧，区分成功、暂时无帧、错误和离线结束\\
pcap\_stats & 读取本句柄的捕获统计，失败不伪造零值\\
pcap\_close & 释放捕获句柄，正常与异常路径均执行\\\bottomrule
\end{tabularx}\end{table}

\subsection{整体流程与退出机制}
\begin{figure}[H]\centering
\begin{tikzpicture}[>=Stealth,node distance=6mm,font=\small,
 box/.style={draw,rounded corners,text width=90mm,align=center,minimum height=9mm}]
\node[box] (args) {参数校验、Npcap 加载与初始化};
\node[box,below=of args] (list) {枚举网卡，交互或参数选择设备};
\node[box,below=of list] (create) {create → 设置参数 → activate};
\node[box,below=of create] (filter) {检查 DLT\_EN10MB，安装可选 BPF，设非阻塞};
\node[box,below=of filter] (check) {检查 Q / Ctrl+C / 数量上限 / 时间上限};
\node[box,below=of check] (next) {pcap\_next\_ex：有帧则边界校验与显示\\无帧则等待 10 ms，错误或离线 EOF 则停止};
\node[box,below=of next] (exit) {读取统计，关闭可选转储、捕获句柄并释放库};
\draw[->] (args)--(list);\draw[->] (list)--(create);\draw[->] (create)--(filter);
\draw[->] (filter)--(check);\draw[->] (check)--(next);\draw[->] (next)--(exit);
\draw[->] (next.west)--++(-.7,0)|-node[pos=.2,left]{继续}(check.west);
\end{tikzpicture}\caption{实时捕获的主要控制流程；离线分支直接打开文件后进入相同解析路径}
\end{figure}

程序不能假定读取缓冲期限会在没有任何流量时定时唤醒阻塞读取，因此主动使用非阻塞模式\cite{nonblock}。
收到状态 1 时立即处理并打印该帧；状态 0 表示暂时没有数据，等待 10 ms 后重新检查停止条件；状态 -1 为错误，-2 表示离线读取结束或 break 情况\cite{nextex}。
驱动返回的头部与数据指针只在后续读取前有效，程序在本轮完成解析和可选转储，不在下一轮继续引用旧指针。

以下是接收分支的简化核心片段，完整实现还记录计数、判断退出条件和处理可选保存：
\begin{code}
const int status = api.pcap_next_ex(handle, &header, &packet);
if (status == 1) {
    print_frame(++received, *header, packet, json);
} else if (status == 0) {
    std::this_thread::sleep_for(std::chrono::milliseconds(10));
} else if (status == PCAP_ERROR_BREAK) {
    break;
} else {
    throw std::runtime_error(api.pcap_geterr(handle));
}
\end{code}

Q 不需要按回车；Ctrl+C 与 Ctrl+Break 通过控制台处理函数仅设置停止标志，资源由主流程释放。
默认最多显示 20 帧或运行 30 秒；两项均可设为零取消限制。时间使用单调时钟，避免系统时钟调整干扰结束判断。

\subsection{输出设计与可选保存}
人类可读主表固定显示 \src{SRC-MAC}、\src{DST-MAC}、\src{Type} 和 \src{After}，并附时间、捕获/原始长度、协议名称。
额外解释单独放在下一行；VLAN 内层解释不会覆盖主表原始字段。
\src{--json} 提供同一组信息的 JSON Lines，供自动化检查，不改变捕获和解析逻辑。

默认仅在屏幕显示，不打开抓包文件。只有明确提供 \src{--save} 时才调用 \src{pcap_dump_open/pcap_dump} 保存原始帧。
\src{--read} 使用 \src{pcap_open_offline} 复核已有文件，输出明确标为“离线复核”，不冒充实时捕获。
设备列表、捕获句柄和可选转储分别由小型 RAII 资源管理器负责释放；Npcap 驱动、DLL 和第三方安装器不打包进提交文件。

\section{实验步骤与测试验证}
\subsection{构建与实时操作}
运行环境先安装 Npcap；重新编译时另需官方 SDK 和 Visual Studio C++ Build Tools。
在仓库根目录执行：
\begin{code}
.\lab02\bootstrap-sdk.ps1
.\lab02\build.ps1
.\lab02\bin\packet_capture.exe --list
.\lab02\bin\packet_capture.exe
\end{code}
最后一条会询问网卡编号。也可显式指定设备、过滤器、数量与时间；以下 6 只是本次枚举编号，不能假定在别的电脑上仍是同一网卡：
\begin{code}
.\lab02\bin\packet_capture.exe --device 6 --count 6 --seconds 30 --filter "icmp"
ping -n 3 10.130.0.1
\end{code}
两条命令在不同终端运行。第二条产生流量；实际演示时应使用现场默认网关，避免照抄本机地址。

\subsection{纯解析器测试}
\ParserTests{} 项 C++ 测试包含：空帧，1--13 字节短帧，恰好 14 字节，MAC 顺序，IPv4/IPv6/ARP 首字节，网络字节序，长度 0 和 1500，保留边界 1501 和 1535，EtherType 起点 1536，单/双 VLAN，截断 VLAN，无负载标签，snaplen 截断，广播 MAC，IPv4 地址、ARP 操作码、全部前缀、重复标签与未对齐缓冲。
这些是明确的构造数据，用于验证边界；真实帧分析单独使用物理网卡捕获。

\subsection{真实 exe 的集成测试}
\begin{longtable}{c p{.68\textwidth}c}
\caption{集成测试覆盖范围与结果}\\\toprule
编号 & 验证内容 & 结果\\\midrule\endfirsthead
\toprule 编号 & 验证内容 & 结果\\\midrule\endhead
01 & Npcap 实际版本、真实设备枚举、IPv4 与物理网卡匹配 & 通过\\
02 & 合成 PCAP：短帧、零首字节、VLAN 原始与内层偏移 & 通过\\
03 & DLT\_NULL 文件明确拒绝，不输出伪 MAC & 通过\\
04 & 非法参数、负数、缺值与不存在的设备 & 通过\\
05 & 非法 BPF 与不存在的输入文件 & 通过\\
06 & 无匹配流量时，仍在 1 秒时限后主动退出 & 通过\\
07 & 实际 ICMP 捕获与 PCAP 原始字段逐条比较 & 通过\\
08 & 同一真实捕获文件在中文路径下读取、过滤 & 通过\\
09 & 数量达到 1 帧后准确退出 & 通过\\
10 & 空闲捕获进程收到 Ctrl+Break 后正常清理退出 & 通过\\
11 & 未指定 save 时不生成或改写抓包文件 & 通过\\\bottomrule
\end{longtable}

记录时间为 \TestTimestamp{}；共 \ParserTests{} 项解析测试和 \IntegrationTests{} 项集成测试全部通过，集成驱动总耗时 \TestDuration{} 秒。
CMake Release 构建和 CTest 解析测试也通过。完整日志与结构化结果分别保存在 evidence 目录的
\src{parser-tests.txt}、\src{integration-tests.txt} 和 \src{results.json}。

\section{真实捕获结果与逐帧分析}\label{sec:analysis}
\subsection{捕获配置与总体结果}
在实际 Wi-Fi 网卡上设置 \src{icmp and host 10.130.0.1}，只捕获与本机网关有关的 ICMP。
默认关闭混杂模式，捕获上限 65535 字节，持续约 \CaptureDuration{} 秒；另一个进程对网关发送三次 32 字节数据的 ping。
共捕获 \PacketCount{} 帧，分别为 3 个 Echo request 和 3 个 Echo reply。
原始 PCAP 为 \src{evidence/real-capture.pcap}，实时程序输出为 \src{evidence/live-capture.jsonl}。

\begin{table}[H]\centering\small\caption{本次真实捕获的统计}
\begin{tabular}{lr}\toprule
指标 & 实测值\\\midrule
应用读取帧数 & \PacketCount\\
短帧 / 截断帧 & \ShortFrames{} / \TruncatedFrames\\
Npcap ps\_recv & \DriverReceived\\
Npcap ps\_drop & \DriverDropped\\
Npcap ps\_ifdrop & \InterfaceDropped\\\bottomrule
\end{tabular}\end{table}

\src{ps_drop} 表示捕获机制因缓冲等原因丢弃的计数；其值为零表示本次句柄没有报告此类捕获丢弃。
\src{ps_ifdrop} 在不同平台可能不受支持，不能把零值当成整个物理网络无丢包的证明\cite{stats}。
三次 ping 均得到应答，是此次请求/应答的观察事实，不能外推为所有网络业务的丢包率。

\subsection{题目必显字段的逐字节核对}
下表从实时 JSON 输出生成；测试另行直接读取 PCAP 记录，比较目的 MAC 的 0--5 字节、源 MAC 的 6--11 字节、12--13 字节的大端值、偏移 14 及捕获/原始长度，六帧全部一致。

\begin{table}[H]\centering\small\caption{真实六帧的必显字段；长度为 caplen/len}
\begin{tabular}{rllllr}\toprule
序号 & 源 MAC & 目的 MAC & 类型 & 首字节 & 长度 \\ \midrule
1 & F8:59:E3:E3:B3:7E & 00:00:5E:00:01:FE & 0x0800 & 0x45 & 74/74 \\
2 & 00:00:5E:00:01:FE & F8:59:E3:E3:B3:7E & 0x0800 & 0x45 & 74/74 \\
3 & F8:59:E3:E3:B3:7E & 00:00:5E:00:01:FE & 0x0800 & 0x45 & 74/74 \\
4 & 00:00:5E:00:01:FE & F8:59:E3:E3:B3:7E & 0x0800 & 0x45 & 74/74 \\
5 & F8:59:E3:E3:B3:7E & 00:00:5E:00:01:FE & 0x0800 & 0x45 & 74/74 \\
6 & 00:00:5E:00:01:FE & F8:59:E3:E3:B3:7E & 0x0800 & 0x45 & 74/74 \\
\bottomrule\end{tabular}

\end{table}

请求方向源 MAC 为 \src{F8:59:E3:E3:B3:7E}、目的 MAC 为 \src{00:00:5E:00:01:FE}；应答方向二者互换。
它们是本次捕获的本地链路地址，不能仅根据目的 IP 就假定所有网络跳点都使用相同 MAC。
类型字段均为 \src{0x0800}，表示后续为 IPv4。字段后首字节均为 \src{0x45}，高四位 4 是版本，低四位 5 表示 IPv4 头长为 $5\times4=20$ 字节。

\subsection{第一帧的十六进制分析}
以下数据直接从第一条真实 PCAP 记录读取，按原始偏移排版，不是手工编造的报文。
\VerbatimInput[fontsize=\small,frame=single,breaklines=true]{first-frame-hex.txt}

偏移 \src{0000} 起六字节为目的 MAC，随后六字节为源 MAC；偏移 12--13 是 \src{08 00}，偏移 14 是 \src{45}。
IPv4 协议字段为 1，表示 ICMP；Echo request 的 ICMP type 为 8、code 为 0。
本次每帧 74 字节，可分解为
\[
14\text{（以太网头）}+20\text{（IPv4头）}+8\text{（ICMP头）}+32\text{（测试数据）}=74.
\]
不能额外把前导码、SFD 或 FCS 加进本次捕获长度；Npcap 提供的是驱动交付的帧视图，FCS 通常不会出现在该缓冲中。

\subsection{请求与应答配对}
根据 IP 方向、ICMP 类型、标识符和序号配对；实际解析结果如下。
\begin{table}[H]\centering\small\caption{六个真实帧的网络层及 ICMP 解释}
\begin{tabular}{rlllr}\toprule
序号 & 源 IPv4 & 目的 IPv4 & ICMP 含义 & 序号字段 \\ \midrule
1 & 10.130.33.54 & 10.130.0.1 & Echo request (8) & 1 \\
2 & 10.130.0.1 & 10.130.33.54 & Echo reply (0) & 1 \\
3 & 10.130.33.54 & 10.130.0.1 & Echo request (8) & 2 \\
4 & 10.130.0.1 & 10.130.33.54 & Echo reply (0) & 2 \\
5 & 10.130.33.54 & 10.130.0.1 & Echo request (8) & 3 \\
6 & 10.130.0.1 & 10.130.33.54 & Echo reply (0) & 3 \\
\bottomrule\end{tabular}

\end{table}
请求的源 IP 为本机 \src{10.130.33.54}，目的为默认网关 \src{10.130.0.1}，应答反向；应答 type 为 0，code 为 0。
逐帧结构化解析另存于 \src{evidence/frame-analysis.json}，其中还包含原始 PCAP 的 SHA256，便于确认表格对应的文件。

\noindent\begin{minipage}{\textwidth}
\subsection{程序界面输出}
以下节选同一份真实 PCAP 在程序中离线复核的前两帧标准输出，明确显示“离线复核”；完整屏幕输出与实时 JSON 记录已保存于 evidence 目录。
这展示了界面的必显列、时间、长度和协议解释。离线读取耗时仅代表读取该文件的时间，不等于上面的真实捕获持续时间。
\VerbatimInput[fontsize=\scriptsize,frame=single,breaklines=true,lastline=7]{capture-screen.txt}
\end{minipage}

\section{问题处理、结论与局限}
\subsection{主要问题及解决}
\begin{enumerate}
\item \textbf{驱动存在却无法加载 API。} 原环境缺少配套 DLL。完成官方 Npcap 安装后，以实际 API 调用和物理网卡捕获验证，而不是仅查驱动文件是否存在。
\item \textbf{回环并非 Ethernet。} 捕获之前检查链路类型；非以太网句柄明确拒绝，避免显示没有意义的“MAC”。
\item \textbf{类型字段与后首字节易混淆。} 固定主表的原始偏移，VLAN 内层字段单独展示；以合成测试验证零值、双标签及截断。
\item \textbf{空闲时阻塞导致无法退出。} 使用非阻塞 next\_ex 与短暂等待，持续检查数量、时间及退出标志；无流量和 Ctrl+Break 测试均通过。
\item \textbf{中文路径受系统代码页影响。} 宽字符入口和 pcap\_init UTF-8 模式共同处理，已用中文文件名验证同一真实捕获可重放。
\item \textbf{屏幕输出需要独立核对。} 将可选保存的真实帧与 JSON 字段逐项比较，避免仅看“程序有输出”就判断解析正确。
\end{enumerate}

\subsection{实验结论}
本实验完成 Npcap 网卡枚举、打开、实时捕获、字段显示与正常退出。使用 C++ 直接处理捕获字节，
\ParserTests{} 项解析测试和 \IntegrationTests{} 项集成测试全部通过；实际六个 ICMP 帧的源/目的 MAC、类型字段、字段后首字节与原始 PCAP 完全一致。
默认只输出屏幕，持久化是可选复核功能；交付含源码、可执行文件和使用本人南开模板完成的 LaTeX 报告。

局限是仅对 DLT\_EN10MB 作以太网字段解析，非以太网类型明确拒绝；实际样本来自一张 Wi-Fi 网卡和一次本机网关测试，VLAN/ARP/IPv6 等边界覆盖不能冒充本次真实网卡观测。
统计只反映该捕获句柄和本次测试，不证明物理网络永不丢包。现场运行仍需安装 Npcap，并按实际网卡编号和网络流量选择设备。

\clearpage
\addcontentsline{toc}{section}{参考文献}
\begin{thebibliography}{9}
\bibitem{npcapdev} Nmap Project. Developing software with Npcap.
\url{https://npcap.com/guide/npcap-devguide.html}.
\bibitem{pcapinit} Nmap Project. pcap\_init.
\url{https://npcap.com/guide/wpcap/pcap_init.html}.
\bibitem{nonblock} Nmap Project. pcap\_setnonblock.
\url{https://npcap.com/guide/wpcap/pcap_setnonblock.html}.
\bibitem{nextex} Nmap Project. pcap\_next\_ex.
\url{https://npcap.com/guide/wpcap/pcap_next_ex.html}.
\bibitem{stats} Nmap Project. pcap\_stats.
\url{https://npcap.com/guide/wpcap/pcap_stats.html}.
\bibitem{ethertypes} IANA. IEEE 802 Numbers.
\url{https://www.iana.org/assignments/ieee-802-numbers/}.
\end{thebibliography}
\end{document}
